\documentclass[11pt]{article}
\usepackage[T1]{fontenc}
\usepackage{lmodern,amsmath,amssymb,microtype}
\usepackage[margin=1in]{geometry}
\usepackage[colorlinks=true,urlcolor=blue]{hyperref}
\hypersetup{pdfauthor={Wentao Li},pdftitle={Shared binary lookups with affine output decoding}}
\newcommand{\F}{\mathbb F_2}
\title{Shared binary lookups with affine output decoding}
\author{Wentao Li\\[2pt]\small Prepared with assistance from OpenAI Codex}
\date{November 2026}
\begin{document}
\maketitle
\begin{abstract}
A lookup on a common address emits a fixed-width binary word. Each output bit is required on a public subset of addresses and is recovered from that word alone. We determine exactly which even subsets of addresses obstruct affine output decoding. Through three specified addresses per output, affine and unrestricted decoders require the same minimum width on every instance. For at most $s$ addresses per output, the worst affine width at conflict chromatic number $k$ is $\lfloor s/2\rfloor\log_2 k+O_s(1)$. This gives the sharp leading cost of requiring an XOR-and-constant output layer.
\end{abstract}

\section{Shared output data}
Heath, Kolesnikov and Ng [1] study garbling an $N$-row, $m$-output lookup table. Here we ask a preprocessing question when each output has a public promise about the rows on which it is used. The encoder $z$ may be an arbitrary table. On a finite nonempty address set $V$ it produces a common word $z(v)\in\F^d$. Output $j$ specifies a function $b_j:D_j\to\F$ on $D_j\subseteq V$ and has a dedicated decoder $g_j$ satisfying $g_j(z(v))=b_j(v)$ there. Outside $D_j$ its value is unconstrained. Each decoder sees only the emitted word. We minimize $d$. We compare arbitrary decoders with decoders of the form $g_j(y)=a_j\cdot y+c_j$ (XORs and constants). The cost is the shared lookup width before garbling; no claim on overall protocol bandwidth is needed.

The conflict graph $G$ joins $u,v$ when some column specifies different bits on them. Every encoding separates conflicting addresses. Conversely, for arbitrary decoders, any proper coloring can be encoded with $\lceil\log_2\chi(G)\rceil$ bits: all addresses of a single color in a column agree there. This is the usual zero-error coloring argument [2].

\section{An exact test for affine decoders}
Let $\mathcal B$ contain the indicators $1_T\in\F^V$ of all even-size subsets $T$ of any $D_j$ for which $\sum_{v\in T}b_j(v)=1$ in $\F$. All these vectors lie in the even-weight subspace $U$ of dimension $|V|-1$. Then an encoding permits all affine decoders exactly when
\[
                   \sum_{v\in T}z(v)\ne0\quad(1_T\in\mathcal B).
\]
Indeed, on a column the equations for $(a_j,c_j)$ have rows $(z(v),1)$. They are inconsistent exactly when some selection of these rows sums to zero with odd sum of right-hand sides. The last coordinate forces the selection to be even. This proves both implications.

Anchor one specified address $v_0$ at $z(v_0)=0$, without cost by translation. Encodings are then in bijection with linear maps $U\to\F^d$: the basis $e_v+e_{v_0}$ for $v\ne v_0$ maps to $z(v)$. Thus the minimum width is the least codimension of a subspace of $U$ avoiding $\mathcal B$. In particular the number of admissible anchored encodings of width $d$ is
\[
 \sum_{J\subseteq\mathcal B}(-1)^{|J|}
           2^{d(|V|-1-\operatorname{rank}J)} .
\]
This follows by inclusion--exclusion on vectors mapped to zero. It is the classical critical problem for binary representable matroids [3], with forbidden vectors explicitly read from the partial table.

\paragraph{Up to three rows.} If $|D_j|\le3$ for all $j$, the only possible nonempty even subsets have size two. The forbidden vectors are exactly the edges of $G$. Thus every proper encoding admits affine decoders, and the minimum is $\lceil\log_2\chi(G)\rceil$. At four rows the equality can fail: take four addresses and a column for every Boolean labeling of the four. Two bits suffice for unrestricted decoding. Four distinct words in $\F^2$ exhaust the space and sum to zero; a column with a single one violates the test. Three bits suffice by encoding the addresses as $0,e_1,e_2,e_3$.

\section{Worst-case width at bounded column size}
Fix $s\ge2$ and $k\ge2$. Let $M_s(k)$ be the maximum minimum affine width over instances with $|D_j|\le s$ and $\chi(G)=k$. Let $D_s(k)$ be the smallest $d$ admitting $k$ labeled points of $\F^d$ on which no nonempty even subset of at most $s$ points sums to zero. This condition includes distinctness via subsets of size two.

\paragraph{Theorem.} $M_s(k)=D_s(k)$. Writing $t=\lfloor s/2\rfloor$, for $k\ge t$,
\[
       \log_2\binom kt\ \le\ M_s(k)\
       \le\ \min\{k-1,\ t\lceil\log_2 k\rceil\}.
\]
In particular $M_s(k)=t\log_2 k+O_s(1)$ as $k\to\infty$.

\paragraph{Exact reduction.}
Take a proper $k$-coloring and encode its colors by the points defining $D_s(k)$. Within a given column repeated colors have consistent bits. Its distinct color points are affinely independent, because an affine dependence over $\F$ is precisely a nonempty even zero sum. Hence the decoder can interpolate their bits. Conversely, take $k$ addresses; for each subset of size at most $s$ include every Boolean labeling as a separate partial column. The conflict graph is complete. Any even zero sum of size at most $s$ conflicts with a column which has a single one on that subset. Thus this instance requires $D_s(k)$, and the maximum is attained.

\paragraph{Bounds.}
The sums of distinct $t$-subsets must be distinct: equality gives a nonempty even zero sum on the symmetric difference, of size at most $2t$. Counting sums gives the lower bound. The points $0,e_1,\ldots,e_{k-1}$ give the first upper bound. For the second, use the binary BCH parity-check construction [4]. Set $m=\lceil\log_2 k\rceil$, choose $k$ distinct $x$ in $\mathbb F_{2^m}$, and encode by
\[
                     x\longmapsto (x,x^3,\ldots,x^{2t-1}).
\]
These are binary words of length $mt$. Suppose $r$ distinct field points with $r$ even, $2\le r\le2t$, had zero vector sum. Their odd power sums through $2t-1$ vanish. Successive squaring gives every positive power sum through $2t-1$; the zeroth sum vanishes by evenness (the constant polynomial has value one also at zero). Consequently the Vandermonde matrix on these $r$ points, with row powers $0,\ldots,r-1$, annihilates the all-ones vector. It is invertible since the points are distinct, a contradiction. Finally $\log_2\binom kt=t\log_2 k-O_t(1)$ for fixed $t$, which finishes the proof. \hfill$\square$

\section*{Sources and status}
[1] D. Heath, V. Kolesnikov and L. K. L. Ng, \href{https://eprint.iacr.org/2024/369}{Garbled Circuit Lookup Tables with Logarithmic Number of Ciphertexts}, EUROCRYPT 2024.\par
[2] H. S. Witsenhausen, \href{https://doi.org/10.1109/TIT.1976.1055607}{The zero-error side information problem and chromatic numbers}, 1976. For characteristic graphs in function computation, see A. Orlitsky and J. R. Roche, \href{https://doi.org/10.1109/18.915643}{Coding for computing}, 2001.\par
[3] H. Crapo and G.-C. Rota, \emph{On the Foundations of Combinatorial Theory: Combinatorial Geometries}, 1970, Ch.~16 (critical problem); see also \href{https://www.sciencedirect.com/science/article/pii/S0195669881800054}{A Constructive Approach to the Critical Problem for Matroids}, 1981.\par
[4] R. C. Bose and D. K. Ray-Chaudhuri, \href{https://www.sciencedirect.com/science/article/pii/S0019995860902874}{On a class of error correcting binary group codes}, 1960.
\par\medskip
The coloring principle, critical-exponent count, and power-sum codes are established mathematics. The proposed contribution is their exact application to shared partial lookups with an affine output layer, including the arity transition. This is a question formulated here. Independent proof and priority review are pending.
\end{document}
